\documentclass[11pt]{article}
\usepackage[margin=1in]{geometry}
\usepackage{amsmath}
\usepackage{amssymb}
\usepackage{hyperref}
\usepackage[dvipsnames]{xcolor}

\DeclareMathOperator*{\argmax}{arg\,max}

\setlength{\parindent}{0pt}
\setlength{\parskip}{1\baselineskip}

\date{2026-04-20}
\author{Benjamin Yeh}
\title{\vspace{-2cm}GRAND-SLAM notes}

\begin{document}

\maketitle
\renewcommand{\baselinestretch}{0.75}\normalsize
\tableofcontents
\renewcommand{\baselinestretch}{1.0}\normalsize

\vfill

\textbf{Disclaimer}: These notes were made solely for the purpose of helping myself understand the GRAND-SLAM model. Relevant parts of the paper were copied or paraphrased directly. I do not seek to infringe on the authors'/journal's copyright or take credit for their work. Much of these notes were written with help from OpenAI Prism and Google Gemini. These notes assume familiarity with statistics and the EM algorithm at the level presented in \href{https://cs229.stanford.edu/main_notes.pdf}{Stanford CS 229} and \href{https://bebi103b.github.io/}{Caltech BE/Bi 103b}.

\textbf{Reference}: Jürges C, Dölken L, Erhard F. Dissecting newly transcribed and old RNA using GRAND-SLAM. \textit{Bioinformatics}. 2018;34(13):i218-i226. doi:\href{https://doi.org/10.1093/bioinformatics/bty256}{10.1093/bioinformatics/bty256}

\pagebreak

\section{Overview}

In the GRAND-SLAM model, each read $i$ has:
\begin{itemize}
    \item $n_i$: the number of thymine positions at which a $T\to C$ mismatch could occur,
    \item $y_i$: the number of observed $T\to C$ mismatches,
    \item $z_i \in \{0,1\}$: a hidden label indicating whether the read came from old RNA ($z_i=0$) or new RNA ($z_i=1$).
\end{itemize}

The hierarchical model is
\[
\begin{gathered}
z_i \sim \mathrm{Bernoulli}(\pi_g) \\
y_i \mid z_i=0; n_i \sim \mathrm{Binomial}(n_i,p_e),
\qquad
y_i \mid z_i=1; n_i \sim \mathrm{Binomial}(n_i,p_c)
\end{gathered}
\]
where \begin{itemize}
    \item $p_e$ (presumably ``e'' for \emph{error}) is the average T to C mismatch rate in unlabeled RNA, and
    \item $p_c$ (presumably ``c'' for \emph{conversion}) is the average T to C mismatch rate in labeled RNA.
\end{itemize}
The probability mass functions are therefore
\[
\begin{gathered}
P(z_i = 0 \mid p_e, p_c, \pi_g) = P(z_i = 0 \mid \pi_g) = 1 - \pi_g, \\
P(z_i = 1 \mid p_e, p_c, \pi_g) = P(z_i = 1 \mid \pi_g) = \pi_g, \\
P(y_i \mid z_i = 0, p_e, p_c, \pi_g; n_i) = P(y_i \mid z_i = 0, p_e; n_i) = \mathrm{Bin}(y_i; n_i, p_e), \\
P(y_i \mid z_i = 1, p_e, p_c, \pi_g; n_i) = P(y_i \mid z_i = 1, p_c; n_i) = \mathrm{Bin}(y_i; n_i, p_c)
\end{gathered}
\]
where
\[
\mathrm{Bin}(y_i;n_i,p) = \binom{n_i}{y_i} p^{y_i} (1 - p)^{n_i - y_i}.
\]

The joint probability mass function can therefore be decomposed as
\[
\begin{aligned}
P(y_i, z_i \mid p_e, p_c, \pi_g; n_i)
&= P(y_i \mid z_i, p_e, p_c, \pi_g; n_i) P(z_i \mid p_e, p_c, \pi_g; n_i) \\
&= P(y_i \mid z_i, p_e, p_c; n_i) P(z_i \mid \pi_g) \\
&= \begin{cases}
    (1 - \pi_g) \mathrm{Bin}(y_i; n_i, p_e), & z_i = 0 \\
    \pi_g \, \mathrm{Bin}(y_i; n_i, p_c), & z_i = 1
\end{cases}
\end{aligned}
\]

The likelihood for the observed data $y_i$ for read $i$ is the two-component binomial mixture (equation 1 from the paper)
\[
\begin{aligned}
P(y_i \mid p_e, p_c, \pi_g; n_i)
&=
\sum_{z_i} P(y_i, z_i \mid p_e, p_c, \pi_g; n_i) \\
&=
\sum_{z_i} P(y_i \mid z_i, p_e, p_c, \pi_g; n_i) P(z_i \mid p_e, p_c, \pi_g; n_i) \\
&= 
(1 - \pi_g)\,\mathrm{Bin}(y_i;n_i,p_e)
+ \pi_g\,\mathrm{Bin}(y_i;n_i,p_c).
\end{aligned}
\]

The probability that a given read is nascent is simply the posterior probability
$$
\begin{aligned}
P(z_i = 1 \mid y_i, p_e, p_c, \pi_g; n_i)
% &= \frac{P(z_i = 1, y_i, p_e, p_c, \pi_g; n_i)}{P(y_i, p_e, p_c, \pi_g; n_i)} \\
% &= \frac{P(y_i \mid z_i = 1, p_e, p_c, \pi_g; n_i) P(z_i = 1 \mid p_e, p_c, \pi_g; n_i) P(p_e, p_c, \pi_g; n_i)}{P(y_i \mid p_e, p_c, \pi_g; n_i) P(p_e, p_c, \pi_g; n_i)} \\
&= \frac{P(y_i \mid z_i = 1, p_e, p_c, \pi_g; n_i) P(z_i = 1 \mid p_e, p_c, \pi_g; n_i)}{P(y_i \mid p_e, p_c, \pi_g; n_i)} \\
&= \frac{P(y_i \mid z_i = 1, p_e, p_c, \pi_g; n_i) P(z_i = 1 \mid p_e, p_c, \pi_g; n_i)}{\sum_{z_i} P(y_i \mid z_i, p_e, p_c, \pi_g; n_i) P(z_i \mid p_e, p_c, \pi_g; n_i)} \\
&= \frac{P(y_i \mid z_i = 1, p_e, p_c; n_i) P(z_i = 1 \mid \pi_g)}{\sum_{z_i} P(y_i \mid z_i, p_e, p_c, \pi_g; n_i) P(z_i \mid \pi_g)} \\
&= \frac{\pi_g\,\mathrm{Bin}(y_i; n_i, p_c)}{(1 - \pi_g)\,\mathrm{Bin}(y_i;n_i,p_e)
+ \pi_g\,\mathrm{Bin}(y_i;n_i,p_c)}
\end{aligned}
$$

The primary objective is to estimate $\pi_g$, the proportion of reads from gene $g$ that are nascent.
$$
f(\pi_g \mid \mathbf{y}, p_e, p_c; \mathbf{n})
= \frac{\left(\prod_i P(y_i \mid p_e, p_c, \pi_g; n_i)\right) P(\pi_g)}{P(\mathbf{y})}
$$
The GRAND-SLAM model uses a uniform prior for $\pi_g$, expressed as a beta distribution with hyperparameters $\alpha = \beta = 1$, yielding
$$
f(\pi_g \mid \mathbf{y}, p_e, p_c; \mathbf{n})
= \frac{\left(\prod_i P(y_i \mid p_e, p_c, \pi_g; n_i)\right) \cdot B(\pi_g \mid \alpha = 1, \beta = 1)}{\int_0^1 \left(\prod_i P(y_i \mid p_e, p_c, \pi_g'; n_i)\right) \cdot B(\pi_g' \mid \alpha = 1, \beta = 1) d\pi_g'}
$$

While the parameters $\pi_g$, $p_e$, and $p_c$ can in theory be estimated together via a single EM procedure, the parameters may be poorly identifiable and the procedure may suffer from poor convergence, as described in the next section. GRAND-SLAM instead fits $p_e$ and $p_c$ separately, then uses fixed estimates of those parameters to estimate $\pi_g$.

The posterior distribution of $\pi_g$ is computed by numerically integrating the evidence (the denominator) via the trapezoidal rule, then approximated via a beta distribution with parameters $\alpha_g$ and $\beta_g$:
$$
f(\pi_g \mid \mathbf{y}, p_e, p_c; \mathbf{n})
\approx b_g(\pi_g; \alpha_g, \beta_g)
= \frac{\pi_g^{\alpha_g - 1} (1 - \pi_g)^{\beta_g - 1}}{B(\alpha_g, \beta_g)}
$$
where $B(\alpha_g, \beta_g) = \frac{\Gamma(\alpha_g) \Gamma(\beta_g)}{\Gamma(\alpha_g + \beta_g)}$ is the beta function.

\section{Parameter Estimation}

\subsection{Standard EM procedure}

\textbf{Step 1: Initialize the parameters.}

Choose initial guesses $p_e^{(0)}$, $p_c^{(0)}$, and $\pi_g^{(0)}$ for every gene $g$.

\textbf{Step 2: E-step.}

For each read $i$ from gene $g$, compute the posterior probability that it came from the \emph{new} RNA component:
\[
\begin{aligned}
r_i^{(t)}
&=
P(z_i=1\mid y_i,p_e^{(t)},p_c^{(t)},\pi_g^{(t)}; n_i) \\
&= \frac{\pi_g^{(t)}\,\mathrm{Bin}(y_i;n_i,p_c^{(t)})}
{(1-\pi_g^{(t)})\,\mathrm{Bin}(y_i;n_i,p_e^{(t)})+\pi_g^{(t)}\,\mathrm{Bin}(y_i;n_i,p_c^{(t)})} \\
&= \frac{\pi_g^{(t)} \binom{n_i}{y_i} \left(p_c^{(t)}\right)^{y_i} \left(1 - p_c^{(t)}\right)^{n_i - y_i}}
{(1 - \pi_g^{(t)}) \binom{n_i}{y_i} \left(p_e^{(t)}\right)^{y_i} \left(1 - p_e^{(t)}\right)^{n_i - y_i} + \pi_g^{(t)} \binom{n_i}{y_i} \left(p_c^{(t)}\right)^{y_i} \left(1 - p_c^{(t)}\right)^{n_i - y_i}} \\
&= \frac{\pi_g^{(t)} \left(p_c^{(t)}\right)^{y_i} \left(1 - p_c^{(t)}\right)^{n_i - y_i}}
{(1 - \pi_g^{(t)}) \left(p_e^{(t)}\right)^{y_i} \left(1 - p_e^{(t)}\right)^{n_i - y_i} + \pi_g^{(t)} \left(p_c^{(t)}\right)^{y_i} \left(1 - p_c^{(t)}\right)^{n_i - y_i}}.
\end{aligned}
\]

Then $1-r_i^{(t)}$ is the posterior probability that the read came from the \emph{old} RNA component.

Intuitively, this step ``softly assigns'' each read to one of the two components.

\textbf{Step 3: M-step for the mixing proportions.}

For each gene $g$, update the estimated fraction of new RNA by averaging those soft assignments over all reads from that gene:
\[
\pi_g^{(t+1)}
=
\frac{1}{m_g}\sum_{i\in g} r_i^{(t)},
\]
where $m_g$ is the number of reads mapped to gene $g$.

\textbf{Step 4: M-step for $p_e$.}

Update the mismatch probability in the old-RNA component by the weighted fraction of mismatches assigned to that component:
\[
p_e^{(t+1)}
=
\frac{\sum_i (1-r_i^{(t)})y_i}
{\sum_i (1-r_i^{(t)})n_i}.
\]

\textbf{Step 5: M-step for $p_c$.}

Update the mismatch probability in the new-RNA component analogously:
\[
p_c^{(t+1)}
=
\frac{\sum_i r_i^{(t)}y_i}
{\sum_i r_i^{(t)}n_i}. 
\]

\textbf{Step 6: Iterate until convergence.}

Repeat the E-step and M-step until the parameter updates are tiny or the log-likelihood stops changing.

This is exactly the usual EM recipe for a mixture model: estimate hidden memberships, then re-estimate the component parameters from those memberships.

\subsubsection{Why the standard EM approach fails here}

Under standard 4sU-labeling conditions, the authors caution that ``$y = 0$ [observing no T to C mismatches] is the mode of both component distributions of the binomial mixture model'' because both $p_e$ and $p_c$ are small. Therefore, the dataset is dominated by observations that are simultaneously very plausible under the old-RNA model and the new-RNA model. For such reads, the E-step responsibility is
\[
r_i^{(t)}
=
\frac{\pi_g^{(t)}(1-p_c^{(t)})^{n_i}}
{(1-\pi_g^{(t)})(1-p_e^{(t)})^{n_i}+\pi_g^{(t)}(1-p_c^{(t)})^{n_i}},
\]
which does not strongly favor either component when both $p_e^{(t)}$ and $p_c^{(t)}$ are small. The likelihood surface therefore has a long, flat ridge rather than a sharp optimum, and EM makes very little progress because its soft assignments barely change from one iteration to the next. This is the practical meaning of the paper's claim that ``standard approaches to estimate pe and pc based on the expectation maximization (EM) algorithm failed.''

GRAND-SLAM's workaround is to first estimate $p_e$ separately, then look only at the tail with sufficiently many mismatches, where the old-RNA component contributes negligibly. In that tail, the data are much more informative about $p_c$, so the estimation problem becomes identifiable enough to solve.

\subsection{\texorpdfstring{Estimating $p_e$}{Estimating pe}}

Methods proposed by Jürges et al.:
\begin{itemize}
    \item Using an additional sample without 4sU labeling, such as a timepoint 0 sample or spike-in RNAs.
    \item Linear regression model to predict T to C error rate from the other error rates (e.g., T to A and T to G error rates).
\end{itemize}

\subsection{\texorpdfstring{Estimating $p_c$}{Estimating pc}}

\subsubsection{Data preprocessing}

The sufficient statistics for parameter estimation are collected in a matrix $A^{(g)}$. Each entry $A_{k,n}^{(g)}$ is the number of reads mapped to a genomic region within gene $g$ containing $n$ thymines with $k$ observed T to C mismatches. Aggregating counts across genes generates the matrix
\[
A=\{A_{k,n}\}=\sum_g A^{(g)}.
\]

Consider reads with a given number of thymine sites $n$ (i.e., a column of $A$). Let $e_{k,n}$ be the number of reads with $k$ mismatches expected to originate from unlabeled RNA. Since $p_e < p_c$ (the underlying assumption of 4sU labeling), there exists a threshold number of T to C mismatches $k$ at and above which only a minor fraction of reads (say, 1\%) originates from the unlabeled component (i.e., $e_{k,n} \leq 0.01 a_{k,n}$). Based on this threshold, we partition the cells $(k, n)$ into 2 mutually exclusive sets:
\begin{itemize}
    \item $R=\{(k,n): e_{k,n} \leq 0.01 a_{k,n}\}$: the inclusion set where unlabeled RNAs are expected to contribute a minor fraction of reads
    \begin{itemize}
        \item Let $R_n=\{k:(k,n)\in R\}$ be the retained mismatch counts for reads with $n$ possible conversion sites.
    \end{itemize}
    \item $X=\{(k,n): e_{k,n}>0.01\,a_{k,n}\}$: the exclusion set where unlabeled RNAs are expected to contribute more than 1\% of reads
    \begin{itemize}
        \item Let $X_n=\{k:(k,n)\in X\}$ be the excluded mismatch counts for reads with $n$ possible conversion sites.
    \end{itemize}
\end{itemize}

How is $e_{k,n}$ defined? Let $M_n = \sum_{k'} A_{k',n}$ be the column sum of $A$: the total number of reads (across all genes) containing $n$ thymine sites. If all of these reads were generated from the unlabeled RNA component ($z_i = 0$), their counts could be collected into a matrix $E = \{E_{k,n}\}$, with the same column sums $M_n$ as in the matrix $A$, but a different distribution within each column. Now, consider a particular value of $n$. Given $p_e$, whether each of $M_n$ reads is observed with $k$ mismatches (and is added to bin $E_{k,n}$) is an independent Bernoulli random variable
$$
\mathbf{1}\{y_i = k\} \sim \mathrm{Bernoulli}(P(y_i = k \mid z_i = 0; n) = \mathrm{Bin}(k; n, p_e))
$$
Therefore, the count of each bin is
$$
E_{k,n} = \sum_{i \, \mid \, n_i = n} \mathbf{1}\{y_i = k\}
$$
with expected value
\[
\begin{aligned}
e_{k,n} = \mathbb{E}[E_{k,n}]
&= \mathbb{E}\left[\sum_{i \, \mid \, n_i = n} \mathbf{1}\{y_i = k\} \right] \\
&= \sum_{i \, \mid \, n_i = n} \mathbb{E}[\mathbf{1}\{y_i = k\}] \\
&= \sum_{i \, \mid \, n_i = n} \mathrm{Bin}(k; n, p_e) \\
&= M_n\,\mathrm{Bin}(k;n,p_e) \\
&= \mathrm{Bin}(k;n,p_e)\sum_{k'} a_{k',n}.
\end{aligned}
\]

\subsubsection{Model: Observed data, latent variables}

Let $A^\text{(unlabeled)}$ and $A^*$ be the unlabeled and labeled counts such that $A^\text{(unlabeled)} + A^* = A$. ($A^\text{(unlabeled)}$ is only introduced here once as a placeholder; it is not used further in the model.)

Consider a latent variable model of the labeled counts $A^*$, where the entries of $A^*$ are partitioned into unobserved and observed sets
\begin{itemize}
    \item $A_\text{unobs}^* = \{a^*_{k,n} \mid (k, n) \in X \}$. For fixed $n$, denote the corresponding subset as $A^*_{\mathrm{unobs},n}$.
    \item $A_\text{obs}^* = \{a^*_{k,n} \mid (k, n) \in R \} = \{ a_{k,n} \mid (k, n) \in R \}$
\end{itemize}
such that $A_\text{unobs}^* \cup A_\text{obs}^* = A^*$ (here we treat $A^* = \{ a_{k,n}^* \}$ as a set rather than a matrix). GRAND-SLAM treats the retained cell counts (defined in the section above) as if they are all generated from the 4sU-labeled component, so the retained cell counts correspond to the observed labeled counts $A_\text{obs}^*$, hence $a^*_{k,n} = a_{k, n}$ for $(k, n) \in R$.

As an example, consider a fixed $n$ (i.e., a column of $A$), say $n = 3$ possible conversion sites. Let's say that the threshold value of $k$ was 2, such that reads with 2 or 3 conversions out of 3 possible conversion sites are now assumed to be from the 4sU-labeled component. (We assume that reads with 0 or 1 conversions may have come from both the unlabeled and 4sU-labeled components.) Therefore, we assume that $A^*_{2,3} = A_{2,3}$ and $A^*_{3,3} = A_{3,3}$ -- these are the observed variables. $A^*_{0, 3}$ and $A^*_{1,3}$ are the unobserved latent variables.

\subsubsection{Complete-data likelihood}

If the complete counts $a^*_{k,n}$ were known, estimating $p_c$ would be easy. The complete-data likelihood is
\[
P(A^* \mid p_c)
=
\prod_n
\left[
\frac{N_n!}{\prod_{k=0}^n a^*_{k,n}!}
\prod_{k=0}^n \mathrm{Bin}(k;n,p_c)^{a^*_{k,n}}
\right].
\]
so the complete-data log-likelihood is
\[
\begin{aligned}
\ell_c(p_c)
&=
\sum_n\sum_{k=0}^n a^*_{k,n}\log \mathrm{Bin}(k;n,p_c)+\text{const} \\
&= \sum_n\sum_{k=0}^n a^*_{k,n} \left(\log\binom{n}{k}+k\log p_c +(n-k)\log(1-p_c) \right) + \text{const} \\
&= \sum_n\sum_{k=0}^n a^*_{k,n} \left(k\log p_c +(n-k)\log(1-p_c) \right) + \text{const} \\
&= \log p_c \left(\sum_n \sum_{k=0}^n a^*_{k,n} k \right) + \log(1-p_c) \left( \sum_n\sum_{k=0}^n a^*_{k,n} (n-k) \right) + \text{const} 
\end{aligned}
\]

Setting the first derivative of the log-likelihood equal to zero, we get
\[
0
= \frac{\partial \ell_c}{\partial p_c}
= \frac{1}{p_c} \left(\sum_n \sum_{k=0}^n a^*_{k,n} k \right)
- \frac{1}{1-p_c} \left( \sum_n\sum_{k=0}^n a^*_{k,n} (n-k) \right).
\]
Rearranging,
\[
\frac{1-p_c}{p_c}
=
\frac{\sum_n\sum_{k=0}^n a^*_{k,n}(n-k)}{\sum_n\sum_{k=0}^n a^*_{k,n}k},
\]
so
\[
p_c
=
\frac{\sum_n\sum_{k=0}^n a^*_{k,n}k}{\sum_n\sum_{k=0}^n a^*_{k,n}k + \sum_n\sum_{k=0}^n a^*_{k,n}(n-k)}.
\]
Since
\[
\sum_n\sum_{k=0}^n a^*_{k,n}k + \sum_n\sum_{k=0}^n a^*_{k,n}(n-k)
=
\sum_n\sum_{k=0}^n a^*_{k,n}n,
\]
this simplifies to
\[
\widehat p_c
=
\frac{\sum_{k,n} k\,a^*_{k,n}}{\sum_{k,n} n\,a^*_{k,n}}.
\]
So with complete data, the solution is just ``total conversions divided by total conversion opportunities.''

\subsubsection{Observed-data likelihood}

Because the excluded cells are unobserved, the actual observed-data likelihood is obtained by summing over all possible unobserved low-mismatch counts. For each fixed $n$, introduce:
\begin{itemize}
    \item $O_n=\sum_{k\in R_n} a_{k,n}$, the total number of retained reads with $n$ possible conversion sites;
    \item $q_n(p_c)=\sum_{k\in R_n}\mathrm{Bin}(k;n,p_c)$, the probability that a read from a 4sU-labeled RNA contains enough mismatches and is retained (observed).
\end{itemize}
Then the retained counts follow a truncated multinomial model
\[
P(A^*_\text{obs} \mid p_c)
=
\prod_n
\left[
\frac{O_n!}{\prod_{k\in R_n} a_{k,n}!}
\prod_{k\in R_n}
\left(\frac{\mathrm{Bin}(k;n,p_c)}{q_n(p_c)}\right)^{a_{k,n}}
\right]
\]
and
\[
\ell_{\mathrm{obs}}(p_c)
=
\sum_n\sum_{k\in R_n} a_{k,n}\log \mathrm{Bin}(k;n,p_c)
-\sum_n O_n\log q_n(p_c)
+\text{const}.
\]
Compared to the complete-data likelihood, this observed-data likelihood is harder to optimize due to the normalizing terms $q_n(p_c)$. EM avoids this difficulty by repeatedly reconstructing the unobserved counts from the current parameter value and then maximizing the simpler complete-data criterion.

We first outline the application of the EM algorithm to this latent variable model, starting with the observed-data log likelihood
$$
\begin{aligned}
\ell_\text{obs}(p_c; A_\text{obs}^*)
&= \log p(A_\text{obs}^* \mid p_c) \\
&= \log \sum_{A_\text{unobs}^*} p(A_\text{obs}^*, A_\text{unobs}^* \mid p_c) \\
&= \log \sum_{A_\text{unobs}^*} Q(A_\text{unobs}^*) \frac{p(A_\text{obs}^*, A_\text{unobs}^* \mid p_c)}{Q(A_\text{unobs}^*)} \\
&\geq \sum_{A_\text{unobs}^*} Q(A_\text{unobs}^*) \log \frac{p(A_\text{obs}^*, A_\text{unobs}^* \mid p_c)}{Q(A_\text{unobs}^*)} \\
&= \mathrm{ELBO}(A_\text{obs}^*, Q, p_c).
\end{aligned}
$$

To make the inequality tight, we set the proposal distribution to the posterior distribution of the unobserved data
$$
Q(A_\text{unobs}^*)
= Q(A_\text{unobs}^* \mid A_\text{obs}^*, p_c^{(t)})
= p(A_\text{unobs}^* \mid A_\text{obs}^*, p_c^{(t)}).
$$

The E-step defines the expected value function
$$
\begin{aligned}
\mathbb{Q}(p_c \mid p_c^{(t)})
&= \mathbb{E}_{A_\text{unobs}^* \sim Q} \log p(A_\text{obs}^*, A_\text{unobs}^* \mid p_c) \\
&= \sum_{A_\text{unobs}^*} Q(A_\text{unobs}^* \mid A_\text{obs}^*, p_c^{(t)}) \log p(A_\text{obs}^*, A_\text{unobs}^* \mid p_c)
\end{aligned}
$$
The M-step then maximizes the expected value function with respect to the parameter $p_c$
$$
p_c^{(t+1)} = \argmax_{p_c} \mathbb{Q}(p_c \mid p_c^{(t)}).
$$

To apply the EM algorithm, we need expressions for the joint log likelihood $\log p(A_\text{obs}^*, A_\text{unobs}^* \mid p_c)$ and the latent variable posterior $p(A_\text{unobs}^* \mid A_\text{obs}^*, p_c^{(t)})$. The expression for the joint log likelihood was already presented earlier in the ``Complete-data likelihood'' section:
\[
\begin{aligned}
\log p(A_\text{obs}^*, A_\text{unobs}^* \mid p_c)
&= \sum_n\sum_{k=0}^n a^*_{k,n}\log \mathrm{Bin}(k;n,p_c) + \text{const} \\
&= \sum_n \left( \sum_{k \in R_n} a_{k,n} \log \mathrm{Bin}(k; n, p_c) + \sum_{k \in X_n} a^*_{k,n} \log \mathrm{Bin}(k; n, p_c) \right) + \text{const}
\end{aligned}
\]

The GRAND-SLAM paper does not specify the specific likelihood model for unobserved counts $A_\text{unobs}^*$; it only presents the E-step and M-step updates. Here, I present one likelihood model suggested by Gemini that is consistent with the GRAND-SLAM E- and M-step updates; other consistent models may exist. To express the posterior distribution for $A_\text{unobs}^*$ explicitly, consider reads from the labeled RNA component with a fixed number of conversion opportunities $n$. Each read is a Bernoulli trial with a probability of ``success'' (retained and observed) $q_n(p_c)$. The number of ``failures'' (reads excluded and unobserved) $U_n = \sum_{k \in X_n} a^*_{k,n}$ before observing $O_n$ ``successes'' therefore follows a Negative Binomial distribution. Given $U_n$, the joint distribution of the specific unobserved counts $\{a^*_{k,n}\}_{k \in X_n}$ scattered across the exclusion bins is a Multinomial distribution with probabilities proportional to $\mathrm{Bin}(k;n,p_c)$. Therefore, for each $n$, $A_{\text{unobs},n}^*$ is drawn from a hierarchical model
$$
\begin{gathered}
U_n \mid A_\text{obs}^*, p_c \sim \mathrm{NB}(O_n, q_n(p_c)) \\
A_{\text{unobs}, n}^* \mid U_n \sim \mathrm{Multinomial} \left(U_n, \left\{ \frac{\mathrm{Bin}(k;n,p_c)}{1-q_n(p_c)} \right\}_{k \in X_n} \right)
\end{gathered}
$$
(More explicitly, $A_{\text{unobs}, n}^* = (a^*_{0,n},\dots,a^*_{k^*,n})$ where $k^* = \max_k \{k \in X_n\}$.) Note that for a given set of unobserved counts $A_{\text{unobs}}^*$, $U_n = \sum_{k \in X_n} a^*_{k,n}$ is fixed. Therefore, the marginal probability of $A_{\text{unobs}}^*$ is simply the joint probability evaluated exactly at $U_n$:
$$
\begin{aligned}
p(A_\text{unobs}^* \mid A_\text{obs}^*, p_c)
&= \prod_n p(A_{\text{unobs},n}^* \mid A_\text{obs}^*, p_c) \\
&= \prod_n p(A_{\text{unobs},n}^*, U_n \mid A_\text{obs}^*, p_c) \\
&= \prod_n p(A_{\text{unobs},n}^* \mid U_n, A_\text{obs}^*, p_c) p(U_n \mid A_\text{obs}^*, p_c) \\
&= \prod_n \left[ \mathrm{Multinomial} \left(A_{\text{unobs},n}^*; U_n, \left\{ \frac{\mathrm{Bin}(k;n,p_c)}{1-q_n(p_c)} \right\}_{k \in X_n} \right) \mathrm{NB}(U_n; O_n, q_n(p_c)) \right]
\end{aligned}
$$

(This expression can be fully expanded\footnote{
$$
\begin{aligned}
p(A_\text{unobs}^* \mid A_\text{obs}^*, p_c)
&= \prod_n \left[ \frac{U_n!}{\prod_{k \in X_n} a^*_{k,n}!} \left(\prod_{k \in X_n} \left( \frac{\mathrm{Bin}(k;n,p_c)}{1 - q_n(p_c)} \right)^{a^*_{k,n}} \right) \binom{O_n + U_n - 1}{U_n} (q_n(p_c))^{O_n} (1 - q_n(p_c))^{U_n} \right] \\
&= \prod_n \left[ \frac{U_n!}{\prod_{k \in X_n} a^*_{k,n}!} \left(\prod_{k \in X_n} \left( \mathrm{Bin}(k;n,p_c) \right)^{a^*_{k,n}}\right) \frac{(O_n + U_n - 1)!}{U_n! (O_n - 1)!} (q_n(p_c))^{O_n} \right] \\
&= \prod_n \left[ \frac{(O_n + U_n - 1)!}{(O_n - 1)! \prod_{k \in X_n} a^*_{k,n}!} (q_n(p_c))^{O_n} \prod_{k \in X_n} \left( \mathrm{Bin}(k;n,p_c) \right)^{a^*_{k,n}} \right]
\end{aligned}
$$}, but it is never used directly.) We are now ready to define the expected value function:
$$
\begin{aligned}
\mathbb{Q}(p_c \mid p_c^{(t)})
&= \mathbb{E}_{A_\text{unobs}^* \sim Q} \log p(A_\text{obs}^*, A_\text{unobs}^* \mid p_c) \\
&= \mathbb{E}_{A_\text{unobs}^* \sim Q} \left[ \sum_n \left( \sum_{k \in R_n} a_{k,n} \log \mathrm{Bin}(k; n, p_c) + \sum_{k \in X_n} a^*_{k,n} \log \mathrm{Bin}(k; n, p_c) \right) + \text{const} \right] \\
&= \sum_n \left( \sum_{k \in R_n} a_{k,n} \log \mathrm{Bin}(k; n, p_c) + \sum_{k \in X_n} \mathbb{E}_{A_\text{unobs}^* \sim Q}[a^*_{k,n}] \log \mathrm{Bin}(k; n, p_c) \right) + \text{const}
\end{aligned}
$$
where the last equality uses the linearity of expectation. This equation shares the exact algebraic form of the complete-data log-likelihood, but with the cell counts replaced by their expected values under $Q$. This naturally defines the expected counts $a^{(t+1)}_{k,n}$:
$$
a^{(t+1)}_{k,n} = \begin{cases} a_{k,n} & k \in R_n \\ \mathbb{E}_Q[a^*_{k,n}] & k \in X_n \end{cases}
$$

To evaluate $\mathbb{E}_Q[a^*_{k,n}]$ for the unobserved bins, use the law of total expectation\footnote{Law of total expectation: $\mathbb{E}[X] = \mathbb{E}[\mathbb{E}[X \mid Y]]$} conditioning on the intermediate variable $U_n$ (the total number of unobserved reads for a given $n$):
$$
\mathbb{E}_Q[a^*_{k,n}] = \mathbb{E}_{U_n \sim Q} \big[ \mathbb{E}[a^*_{k,n} \mid U_n] \big]
$$
The inner conditional expectation is the mean of a single component in a Multinomial distribution, which is the total number of trials $U_n$ multiplied by the probability of falling into category $k$:
$$
\mathbb{E}[a^*_{k,n} \mid U_n] = U_n \frac{\mathrm{Bin}(k;n,p_c^{(t)})}{1-q_n(p_c^{(t)})}
$$
Substitute the inner expectation back into the outer expectation. Because $U_n$ is the only random variable, factor out the probabilities:
$$
\mathbb{E}_Q[a^*_{k,n}] = \mathbb{E}_{U_n} [ U_n ] \left( \frac{\mathrm{Bin}(k;n,p_c^{(t)})}{1-q_n(p_c^{(t)})} \right)
$$
The expectation $\mathbb{E}_{U_n} [ U_n ]$ is the mean of a Negative Binomial distribution representing the number of failures before $O_n$ successes, with a success probability of $q_n(p_c^{(t)})$. The mean of this distribution is:
$$
\mathbb{E}_{U_n}[U_n] = O_n \frac{1 - q_n(p_c^{(t)})}{q_n(p_c^{(t)})}
$$
Thus, we have
$$
\begin{aligned}
\mathbb{E}_Q[a^*_{k,n}]
&= \left( O_n \frac{1 - q_n(p_c^{(t)})}{q_n(p_c^{(t)})} \right) \left( \frac{\mathrm{Bin}(k;n,p_c^{(t)})}{1-q_n(p_c^{(t)})} \right) \\
&= \frac{O_n \mathrm{Bin}(k;n,p_c^{(t)})}{q_n(p_c^{(t)})} \\
&= \frac{\left(\sum_{k' \in R_n} a_{k',n}\right) \mathrm{Bin}(k;n,p_c^{(t)})}{\sum_{k'\in R_n}\mathrm{Bin}(k';n,p_c^{(t)})}
\end{aligned}
$$
This finally gives equation (5) from the GRAND-SLAM paper,
$$
a^{(t+1)}_{k,n} = \begin{cases}
a_{k,n} & k \in R_n \\
\frac{\left(\sum_{k' \in R_n} a_{k',n}\right) \mathrm{Bin}(k;n,p_c^{(t)})}{\sum_{k'\in R_n}\mathrm{Bin}(k';n,p_c^{(t)})} & k \in X_n
\end{cases}
$$
which we plug into the expected value function
$$
\begin{aligned}
\mathbb{Q}(p_c \mid p_c^{(t)})
&= \mathbb{E}_{A_\text{unobs}^* \sim Q} \log p(A_\text{obs}^*, A_\text{unobs}^* \mid p_c) \\
&= \sum_n \sum_{k=0}^n a^{(t+1)}_{k,n} \log \mathrm{Bin}(k;n,p_c)+\text{const}.
\end{aligned}
$$
In the M-step, we maximize $\mathbb{Q}(p_c \mid p_c^{(t)})$ over $p_c$. Since the form of $\mathbb{Q}(p_c \mid p_c^{(t)})$ is identical to the form of the complete-data log-likelihood $\ell_c(p_c)$, its maximum is of the same form as that of the complete-data log-likelihood
\[
p_c^{(t+1)}
=
\frac{\sum_{k,n} k\,a^{(t+1)}_{k,n}}{\sum_{k,n} n\,a^{(t+1)}_{k,n}}.
\]

\subsubsection{GRAND-SLAM's bisection scheme}

The GRAND-SLAM paper writes, ``We noticed that running the EM algorithm led to extremely slow convergence rates. Thus, we use the following bisection scheme instead: For the search interval $[l,r]$ (starting with $l=0$ and $r=1$), we set $p^{(t)} = \frac{l+r}{2}$ and compute $p^{(t+1)}$ by a single EM iteration. If $p^{(t+1)} < p^{(t)}$, we continue with the search interval $[l, p^{(t)}]$, otherwise with $[p^{(t)}, r]$. We stop if $r-l < 10^{-8}$.''

Properties of this bisection scheme:\begin{itemize}
\item Guaranteed to converge within an interval of $10^{-8}$ in exactly $\lceil \log_2(10^8) \rceil = 27$ iterations.
\item Does not share the monotonically-improving likelihood guarantee of the standard EM algorithm. \textcolor{red}{(need proof) However, since the log-likelihood function is concave and unimodal, the bisection scheme is still guaranteed to converge to the global maximum $\hat{p}$.}
\item Does not assume that $p^{(t+1)}$ falls in the interval $[l, r]$; the next sub-interval is chosen solely based on the boolean condition $p^{(t+1)} < p^{(t)}$. \textcolor{red}{(need proof) Nonetheless, for the GRAND-SLAM model, $p^{(t+1)}$ will always fall inside $[l, r]$, since for a concave log-likelihood of a scalar parameter, a single EM step will never overshoot the maximum likelihood estimate.}
\end{itemize}

\section{Estimating RNA half-life}

State variable
\begin{itemize}
    \item $a$: total abundance of RNA species
\end{itemize}
Parameters
\begin{itemize}
    \item $\sigma$: transcription rate
    \item $\delta$: decay rate
\end{itemize}
The governing first-order ODE model is
$$
\frac{da}{dt} = \sigma - \delta a(t), \qquad a(0) = a_0
$$
with solution to the initial value problem
$$
a(t) = \left(a_0 - \frac{\sigma}{\delta} \right) e^{-t \delta} + \frac{\sigma}{\delta}
$$
and steady-state $\frac{\sigma}{\delta}$. This ODE model is used to separately describe the abundance of pre-existing RNA $a_\text{pre}$ and newly synthesized RNA $a_\text{new}$.
\begin{itemize}
\item Pre-existing RNA is assumed to start with steady-state abundance and then decay without a production term. \textbf{This corresponds to the physical assumption that all new RNAs synthesized after the start of 4sU pulse are in the 4sU-labeled component with T to C conversion rate $p_c$.}
  $$
  a_\text{pre}(t) = \frac{\sigma}{\delta} e^{-t \delta}
  $$
\item New RNA starts from an initial abundance of 0.
  $$
  a_\text{new}(t) = - \frac{\sigma}{\delta} e^{-t \delta} + \frac{\sigma}{\delta}
  $$
\end{itemize}
The total abundance is maintained at steady-state, $a_\text{pre}(t) + a_\text{new}(t) = \frac{\sigma}{\delta}$. At any time $t$, the proportion of new RNA out of total RNA is
$$
\begin{aligned}
\pi(t)
&= \frac{a_\text{new}(t)}{a_\text{new}(t) + a_\text{pre}(t)} \\
&= \frac{- \frac{\sigma}{\delta} e^{-t \delta} + \frac{\sigma}{\delta}}{- \frac{\sigma}{\delta} e^{-t \delta} + \frac{\sigma}{\delta} + \frac{\sigma}{\delta} e^{-t \delta}} \\
&= 1 - e^{-t \delta}
\end{aligned}
$$
This can be used to estimate the decay rate (here, expressing the decay rate $\delta$ as a function of the random variable $\pi$, for which we have approximated its posterior distribution with a beta distribution $b_g(\pi_g; \alpha_g, \beta_g)$)
$$
\delta_t(\pi) = -\frac{1}{t} \log(1 - \pi)
$$
with inverse function
$$
\pi_t(\delta) = 1 - e^{-t \delta}.
$$
Then, the posterior density function of $\delta_t$ (posterior density because we are using the posterior density of $\pi_g$) is
$$
\begin{aligned}
d(\delta; \alpha_g, \beta_g)
&= \beta_g(\pi_t(\delta); \alpha_g, \beta_g) \left\lvert \frac{d\pi_t(\delta)}{d\delta} \right\rvert \\
&= \frac{(\pi_t(\delta))^{\alpha_g - 1} (1 - (\pi_t(\delta)))^{\beta_g - 1}}{B(\alpha_g, \beta_g)} t e^{-t \delta} \\
&= \frac{(1 - e^{-t \delta})^{\alpha_g - 1} (e^{-t \delta})^{\beta_g - 1}}{B(\alpha_g, \beta_g)} t e^{-t \delta} \\
&= \frac{t}{B(\alpha_g, \beta_g)} (1 - e^{-t \delta})^{\alpha_g - 1} e^{-t \beta_g \delta + t \delta} e^{-t \delta} \\
&= \frac{t}{B(\alpha_g, \beta_g)} (1 - e^{-t \delta})^{\alpha_g - 1} e^{-t \beta_g \delta}
\end{aligned}
$$
with log posterior density
$$
\ell(\delta) = (\alpha_g - 1) \log(1 - e^{-t \delta}) - t \beta_g \delta + \text{const}.
$$

Generalized to multiple observations across differnt timepoints, if approximate posterior beta densities with parameters $(\alpha_i, \beta_i)$ are available at labeling times $t_1, \dots, t_m$, then—assuming the experiments are conditionally independent given a shared decay rate—the combined log-posterior for $\delta$ is, up to an additive constant,
$$
\ell(\delta)
= \sum_{i=1}^m (\alpha_i - 1) \log\left(1 - e^{-t_i \delta}\right) - t_i \beta_i \delta + \text{const},
$$
which is equation (17) in Jürges et al. The MAP estimate is found numerically:
$$
\hat{\delta}_{\mathrm{MAP}} = \argmax_{\delta \ge 0} \sum_{i=1}^m (\alpha_i - 1) \log\left(1 - e^{-t_i \delta}\right) - t_i \beta_i \delta.
$$
and can be converted to the RNA half-life
$$
\hat{k} = \frac{\log 2}{\hat{\delta}_{\mathrm{MAP}}}.
$$

\end{document}